%==============================================================================
\section{Limitations --- Giới hạn của học máy}
%==============================================================================

\begin{dongco}[title={ML mạnh, nhưng không phải lúc nào cũng phù hợp}]
Học máy là công nghệ mạnh và đã thay đổi cách ta phân tích dữ liệu.
Tuy nhiên, như mọi công nghệ, ML có những giới hạn cần hiểu rõ để dùng đúng.
\end{dongco}

\subsection{Dependence on Data Quality}

\begin{dinhnghia}[title={Phụ thuộc vào chất lượng dữ liệu}]
Mô hình ML “tốt đến đâu” phụ thuộc vào dữ liệu huấn luyện.
Nếu dữ liệu thiếu, thiên lệch hoặc chất lượng kém, mô hình có thể hoạt động không tốt.
\end{dinhnghia}

\subsection{Lack of Transparency}

\begin{dinhnghia}[title={Thiếu minh bạch}]
Mô hình ML có thể rất phức tạp, khiến khó hiểu mô hình đi đến dự đoán như thế nào.
Thiếu minh bạch làm khó việc giải thích kết quả với stakeholder.
\end{dinhnghia}

\subsection{Limited Applicability}

\begin{dinhnghia}[title={Khả năng áp dụng có giới hạn}]
Mô hình ML được thiết kế để tìm pattern trong dữ liệu, nên không phù hợp với mọi loại dữ liệu hay mọi loại bài toán.
\end{dinhnghia}

\subsection{High Computational Costs}

\begin{dinhnghia}[title={Chi phí tính toán cao}]
Mô hình ML có thể tốn tài nguyên tính toán, cần nhiều năng lực xử lý và lưu trữ.
\end{dinhnghia}

\subsection{Data Privacy Concerns}

\begin{dinhnghia}[title={Lo ngại về quyền riêng tư dữ liệu}]
Mô hình ML đôi khi thu thập và dùng dữ liệu cá nhân, dẫn tới lo ngại về quyền riêng tư và an toàn dữ liệu.
\end{dinhnghia}

\subsection{Ethical Considerations}

\begin{dinhnghia}[title={Vấn đề đạo đức}]
Mô hình ML có thể duy trì/lan truyền thiên lệch hoặc phân biệt đối xử với một số nhóm, tạo ra vấn đề đạo đức.
\end{dinhnghia}

\subsection{Dependence on Experts}

\begin{dinhnghia}[title={Phụ thuộc vào chuyên gia}]
Phát triển và triển khai mô hình ML cần chuyên môn về data science, thống kê và lập trình.
Vì vậy, với tổ chức không có sẵn kỹ năng này, việc triển khai ML có thể khó khăn.
\end{dinhnghia}

\subsection{Lack of Creativity and Intuition}

\begin{dinhnghia}[title={Thiếu sáng tạo và trực giác}]
Thuật toán ML giỏi tìm pattern trong dữ liệu nhưng thiếu sáng tạo và trực giác.
Vì vậy, ML có thể không giải tốt bài toán cần tư duy sáng tạo hoặc trực giác.
\end{dinhnghia}

\subsection{Limited Interpretability}

\begin{dinhnghia}[title={Khả năng diễn giải hạn chế}]
Một số mô hình (ví dụ deep neural networks) khó diễn giải.
Điều này làm khó hiểu vì sao mô hình đưa ra dự đoán.
\end{dinhnghia}